\section{Mathematical Reasoning}

\answer{AG2-1}
Removing the known diagonal's contribution at vertices $1,3$ leaves zero residual parity at every vertex. The two incident cycle bits must agree at each vertex, so propagation around the cycle makes all four bits zero or all four one. There are exactly two solutions. Four independent fair bits have $16$ equally likely assignments, giving probability $2/16=1/8$. Observing $12=0$ leaves only the all-zero solution; the remaining unknown graph is a path, consistent with A19's forest criterion.

\answer{AG2-2}
The column and two diagonal slopes are $a,a+1,1-a$, all of which must be nonzero modulo $5$. Excluding $a=0,1,4$ leaves $a=2,3$. Each permits all five shifts $b$, for $10$ arrangements. Since $9$ is divisible by $3$, the square-sum necessary condition in A14 fails. No arrangement exists, and using a non-affine rule cannot avoid that obstruction.

\answer{AG2-3}
The upper endpoint of the jump is $3/8+1/4=5/8$, so the conditional law is $\operatorname{Unif}(3/8,5/8)$. Half of that interval lies above $1/2$, giving joint probability $(1/4)(1/2)=1/8$.
The original variable takes values in the countable set $\{r_j\}$, so $h(X)$ almost surely belongs to the countable image $\{h(r_j)\}$. A continuous uniform law assigns probability zero to every countable set, contradicting probability one here. Density of the atom locations on the real line does not change this argument.

\answer{AG2-4}
The deleted member held three types, each of which changes from three holders to two. The other four types still have three holders. Thus there are $3$ two-holder and $4$ three-holder types, and none disappears. No item held by a remaining member is deleted, so each still has three. The unique item shared by each remaining pair is also retained.
Independently, let the two counts be $x,y$. Person--item incidences and shared pairs give $2x+3y=6\cdot3=18$ and $x+3y=\binom62=15$. Subtracting yields $x=3$, then $y=4$, agreeing with the construction.

\answer{AG3-1}
Uniform hiding makes every search succeed with probability $1/2$, while uniform choice of search pair covers every location with probability $1/2$. Hence the value is $1/2$.
For hiding probabilities $p_i$, the average success probability of the four searches is $1/2$. Their maximum can be at most $1/2$ only if all are equal. Thus all optimal hiding strategies are $p_1=p_3=a$, $p_2=p_4=1/2-a$, where $0\le a\le1/2$.
For the searcher, the average coverage of the four locations is also $1/2$. Minimum coverage is at least $1/2$ exactly when every location has that coverage. Writing $q_{ij}$ for the probability of searching edge $ij$, all solutions are
\[
 q_{12}=q_{34}=b,\qquad q_{23}=q_{41}=1/2-b,
 \qquad0\le b\le1/2.
\]
The conditions are necessary and sufficient, including the interval endpoints.

\answer{AG3-2}
Pile sizes $1,2,3$ are respectively losing, winning, losing: one stone must be taken at a loss; two can leave one; from three, taking one leaves a winning state and taking all loses immediately. From $4,5,6,7$, removing $3,4,3,4$ respectively leaves $1$ or $3$. Thus only $1,3$ lose among the first seven positive positions.
The losing residues are $1,3\pmod7$. Removing $1,3,4$ from either cannot reach another losing residue; a move that empties the pile is itself losing. From the other residues $0,2,4,5,6$, remove respectively $4,1,3,4,3$ to leave a positive losing position. These moves are legal even at the smallest positive representatives. Strong induction proves the classification.
Since $20\equiv6\pmod7$, remove $3$ to leave $17\equiv3$. The other moves leave winning states $19,16$, so this is the unique winning first move. Under A09's normal-play rule, remove $4$ to leave $16\equiv2$; that first move is also uniquely winning.

\answer{AG3-3}
With both servers occupied, no-deviation conditions become
\[
 k\le\frac{9-k}{2}+\delta,\qquad
 \frac{8-k}{2}+\delta\le k+1,
\]
equivalently $6+2\delta\le3k\le9+2\delta$. For $k=3$ to be an equilibrium, one needs $0\le\delta\le3/2$. In this range neither empty-server configuration is an equilibrium. For $0<\delta<3/2$, the only integer solution is $k=3$.
At $\delta=0$, occupancy $k=2$ is also an equilibrium; at $\delta=3/2$, so is $k=4$. Outside the range, $k=3$ itself violates a no-deviation condition. Pure equilibria always exist: for $0\le\delta<15/2$, take $k=\lceil2+2\delta/3\rceil$; for $\delta\ge15/2$, all eight users choosing the slow server is an equilibrium. Thus the desired fees are exactly $0<\delta<3/2$. Total service cost at that occupancy remains $43/2$: fees change incentives, not the actual service cost of a fixed occupancy.

\answer{AG3-4}
Customers split at $(a+b)/2$. Integrating over the resulting intervals gives
\[
 C(a,b)=\frac{a^2}{2}+\frac{(b-a)^2}{4}
       +\frac{(1-b)^2}{2}.
\]
The gradient is $((3a-b)/2,(3b-a-2)/2)$, with unique stationary point $a=1/4,b=3/4$. The Hessian has eigenvalues $1,2$, so the objective is strictly convex. This feasible stationary point is the unique ordered global optimum, and substitution gives $C=1/8$.
Under competition, separated providers cannot be in equilibrium: the left one can move toward the right one and increase its share. Coincidence must be at the uniform median $1/2$; elsewhere, a small deviation toward the side containing more demand earns more than half. At the midpoint, no deviation improves on half. Thus the unique pure equilibrium is $(1/2,1/2)$, with mean distance $\int_0^1|t-1/2|\,\dd t=1/4$, twice the cooperative minimum.

\par\Needspace{8\baselineskip}
\subsection*{Section Exercises}
\answer{OA17}
Suppose there are $N$ types in total, and type $j$ occurs in $r_j$ collections, where $1\le r_j\le4$. Counting all memberships gives
\[
 \sum_{j=1}^N r_j=40.
\]
Now add the sizes of the six pairwise intersections. Type $j$ is counted $\binom{r_j}{2}$ times, so
\[
 \sum_{j=1}^N\binom{r_j}{2}\le36,
 \qquad \sum_{j=1}^N r_j^2\le112.
\]
By Cauchy--Schwarz,
\[
 40^2\le N\sum_{j=1}^N r_j^2\le112N,
 \qquad N\ge\frac{100}{7}.
\]
Since $N$ is an integer, $N\ge15$.

Label fifteen types by $1$ through $15$ and take
\[
 \begin{aligned}
 A&=\{1,2,3,4,5,6,7,8,9,10\},\\
 B&=\{1,2,3,4,5,6,11,12,13,14\},\\
 C&=\{1,2,7,8,9,10,11,12,13,14\},\\
 D&=\{3,4,5,7,8,9,11,12,13,15\}.
 \end{aligned}
\]
Each collection has ten elements, every pairwise intersection has exactly six elements, and their union contains all fifteen types. The minimum is therefore $\boxed{15}$.

\answer{OA22}
If $x>y$, then $3x-2y=x+2(x-y)\ge4$. If $x=y$, integrality excludes $x=1$, so $3x-2y=x\ge2$. It remains to consider $y=x+d$ with $d\ge1$.

Set
\[
 P=(x+1)(x+2)\cdots(x+d).
\]
The given fraction is $x!P/(P+1)$. Since $P$ and $P+1$ are coprime, integrality is equivalent to $P+1\mid x!$.

Suppose, for a contradiction, that $x-2d\le1$, or $x\le2d+1$. Let $p$ be any prime divisor of $P+1$. Divisibility gives $p\le x$. If $p\le d$, one of the $d$ consecutive integers $x+1,\ldots,x+d$ is divisible by $p$, contradicting $p\nmid P$. Hence $p>d$. If also $p\le y/2$, then
\[
 x<2d+2\le2p\le y.
\]
Thus $2p$ is a factor in the product $P$, again a contradiction. Every prime divisor must therefore lie in
\[
 \frac y2<p\le x.
\]
If $d\ge x$, this interval is empty, contradicting $P+1>1$. Otherwise the number of integers in the interval is
\[
 x-\left\lfloor\frac{x+d}{2}\right\rfloor
 =\left\lceil\frac{x-d}{2}\right\rceil\le d,
\]
where the last bound follows from $x\le2d+1$ and $d\ge1$. Moreover, $p>y/2>x/2$, so each such prime occurs only once in $x!$. Since $P+1\mid x!$, the integer $P+1$ is a product of at most $d$ distinct primes, each at most $x$. Consequently,
\[
 P+1\le x^d<(x+1)\cdots(x+d)=P,
\]
a contradiction. Hence $3x-2y\ge2$. For $(x,y)=(4,5)$, the given fraction equals $20$ and $3x-2y=2$, so the constant is sharp.

\answer{OA28}
The total weight is $\sum_{n=0}^\infty2^{-n}=2$. Thus the series converges at every $x$, with a tail bound that is uniform in $x$.

First let $x$ be irrational. Given $\varepsilon>0$, choose $N$ such that $\sum_{n>N}2^{-n}<\varepsilon$. None of the finitely many points $r_0,\ldots,r_N$ equals $x$. Hence some $\delta>0$ ensures that, whenever $|y-x|<\delta$, each of these rational points is on the same side of $y$ as it is of $x$. The first $N+1$ terms remain unchanged, so
\[
 |f(y)-f(x)|\le\sum_{n>N}2^{-n}<\varepsilon.
\]
The function is continuous at every irrational point.

Now take $x=r_k$ and separate its term:
\[
 f(y)=2^{-k}\mathbf1_{\{y>r_k\}}+g(y).
\]
The same finite-prefix and tail argument shows that $g$ is continuous at $r_k$, because its prefix contains no point equal to $r_k$. The strict inequality in the definition gives $f(r_k)=g(r_k)$. Therefore
\[
 \boxed{\lim_{y\uparrow r_k}f(y)=f(r_k),\qquad
 \lim_{y\downarrow r_k}f(y)=f(r_k)+2^{-k}.}
\]
Each rational point has an upward jump of size $2^{-k}$ and is left-continuous but not continuous. The continuity points are exactly the irrationals. Dense jump locations cause no contradiction: their tail weights can be made uniformly small.

\answer{OA09}
For a positive integer $r$, let $\ell(r)$ be the largest power of two dividing $r$, equivalently the value of its least significant nonzero binary digit. Consider a general position $(r,m)$ with $r$ stones remaining and a current limit of $m$. We prove
\[
 (r,m)\text{ is losing}\quad\Longleftrightarrow\quad m<\ell(r).
\]
Use induction on the number of stones. If $m\ge\ell(r)$, take $\ell(r)$ stones. This either wins immediately or leaves $r-\ell(r)$, whose least significant nonzero bit is strictly larger than $\ell(r)$. The next limit is $\ell(r)$, so the opponent receives a losing position.

If $m<\ell(r)$, every legal move $a$ satisfies $1\le a<\ell(r)$. Write $a=2^s q$ with $q$ odd. Because $r$ is divisible by a higher power of two, $r-a$ has exactly $s$ factors of two. Thus
\[
 \ell(r-a)=\ell(a)\le a.
\]
The opponent's limit is $a$, so the resulting position is winning. Every legal move therefore gives a winning position to the opponent, completing the induction.

Initially the position is $(n,n-1)$. It is losing exactly when $\ell(n)=n$, which means that $n$ is a power of two. Hence
\[
 \boxed{\text{First player wins}\iff n\notin\{2^k:k\ge1\}.}
\]
In a winning initial position, take $\ell(n)$ stones. On each subsequent turn, take the least significant nonzero bit of the remaining count. The position criterion ensures that this move is legal and eventually takes the last stone.

\answer{OA13}
The first player chooses $6$. Every later legal number must avoid the prime factors $2$ and $3$, leaving only
\[
 5,7,11,13,17,19,23,29
\]
as possible prime factors. The product of any two distinct primes in this list is at least $5\cdot7=35>30$. Thus every later legal number uses exactly one prime from the list; for example, $25$ uses only the prime $5$. Once a number has been selected, its prime factor cannot be used again.

Each move therefore consumes exactly one of the eight available primes. Conversely, if a prime $p$ has not yet been used, the number $p$ itself remains a legal choice. The game cannot end while such a prime remains. There are consequently exactly eight further moves after the initial $6$, regardless of the players' later choices.

There are nine moves in total, so the first player makes the last move and wins. Both parts of the argument matter: every move consumes one prime, and every unused prime guarantees another legal move.

